% =============================================================
%  Module Description — Linear Algebra (M103-S27)
%  Shown on https://wosm.academy as the course page for M103-S27.
%  Compile with: pdflatex m103-module-description.tex
%  Write "TODO(owner): ..." where information is still missing;
%  the website highlights it.
% =============================================================
\documentclass[11pt,a4paper]{article}

\usepackage[utf8]{inputenc}
\usepackage[T1]{fontenc}
\usepackage{lmodern}
\usepackage[margin=2.2cm]{geometry}
\usepackage{amssymb}
\usepackage{booktabs}
\usepackage{tabularx}
\usepackage{array}
\usepackage{enumitem}
\usepackage{xcolor}
\usepackage{titlesec}
\usepackage{fancyhdr}
\usepackage[hidelinks]{hyperref}

% ---------- Module details (edit these) ----------------------
% \ModuleCode must equal the course's `code` in src/content/courses/.
\newcommand{\ModuleTitle}{Linear Algebra}
\newcommand{\ModuleCode}{M103-S27}
\newcommand{\Programme}{Foundation programme}
\newcommand{\Level}{Semester 2}
\newcommand{\Semester}{Spring 2027}
\newcommand{\Credits}{7 ECTS}
\newcommand{\ModuleLeader}{Dr. Ali Khudiyev}
\newcommand{\Email}{ali.khudiyev@wosm.academy}
\newcommand{\OfficeHours}{Monday--Friday, 10:00--17:00 (Baku time, UTC+4)}
\newcommand{\Prerequisites}{Discrete Structures (M101-F26)}

% ---------- Styling ------------------------------------------
\definecolor{accent}{RGB}{0,70,130}
\titleformat{\section}{\large\bfseries\color{accent}}{\thesection.}{0.5em}{}
\titlespacing*{\section}{0pt}{1.2em}{0.5em}
\setlist{itemsep=2pt, topsep=4pt}
\setlength{\parindent}{0pt}
\setlength{\parskip}{4pt}

\pagestyle{fancy}
\fancyhf{}
\lhead{\small \ModuleCode\ -- \ModuleTitle}
\rhead{\small Module Description}
\cfoot{\small \thepage}
\renewcommand{\headrulewidth}{0.4pt}

% =============================================================
\begin{document}

\begin{center}
  {\LARGE\bfseries\color{accent} \ModuleTitle}\\[4pt]
  {\large Module Description}
\end{center}
\vspace{6pt}

% ---------- Summary table (not shown on the website: the course page
% shows these details from the macros above) --------------------------
\begin{tabularx}{\textwidth}{>{\bfseries}l X >{\bfseries}l X}
\toprule
Module code   & \ModuleCode    & Credits      & \Credits \\
Programme     & \Programme     & Level        & \Level \\
Semester      & \Semester      & Prerequisites & \Prerequisites \\
Module leader & \ModuleLeader  & Email        & \href{mailto:\Email}{\Email} \\
Office hours  & \multicolumn{3}{l}{\OfficeHours} \\
\bottomrule
\end{tabularx}

% ---------- Sections (shown on the website) ------------------
\section{Module Overview}
Linear Algebra is the study of vectors, linear maps and the matrices that
represent them. It is one of the most widely used areas of mathematics in computer
science: graphics, machine learning, data analysis, search engines, coding theory
and scientific computing all rest on it.

The module covers:
\begin{itemize}
  \item Systems of linear equations, Gaussian elimination and matrix algebra
        (first 4--5 weeks)
  \item Vector spaces, linear maps, bases and dimension, determinants
  \item Inner products, orthogonality and least squares
  \item Eigenvalues, diagonalisation, the spectral theorem and the singular value
        decomposition
\end{itemize}

The module balances the abstract view (vector spaces and linear maps) with
computation (matrices and algorithms). Students implement the core algorithms in
short programming exercises and meet applications such as PageRank, image
compression and linear regression. It builds on the sets, functions and algebraic
structures of Discrete Structures and is required for Probability Theory and
Statistics.

\section{Aims}
\begin{itemize}
  \item Introduce vector spaces and linear maps as the common language for many problems in computing and science.
  \item Develop fluency with matrix computations and the algorithms behind them, including their cost and numerical behaviour.
  \item Build geometric intuition for subspaces, projections, orthogonality and eigenvectors.
  \item Prepare students for probability, statistics, machine learning, graphics and numerical methods.
\end{itemize}

\section{Learning Outcomes}
On successful completion of this module, students will be able to:
\begin{enumerate}[label=LO\arabic*.]
  \item Solve systems of linear equations with Gaussian elimination and LU factorisation, and describe the solution set in terms of rank and null space.
  \item Work with abstract vector spaces: decide linear independence, find bases and dimensions, and apply the rank--nullity theorem.
  \item Represent linear maps by matrices, change basis, and compute and interpret determinants.
  \item Use inner products to compute orthogonal projections, apply Gram--Schmidt and QR factorisation, and solve least-squares problems.
  \item Compute eigenvalues and eigenvectors, diagonalise matrices where possible, and apply the spectral theorem to symmetric matrices.
  \item Compute and interpret the singular value decomposition and apply it to low-rank approximation and data analysis.
  \item Implement core linear algebra algorithms in code and estimate their computational cost.
\end{enumerate}

\section{Syllabus}
\begin{tabularx}{\textwidth}{c X}
\toprule
\textbf{Unit} & \textbf{Topic} \\
\midrule
1  & Vectors in $\mathbb{R}^n$: linear combinations, dot product, length, angle \\
2  & Systems of linear equations; Gaussian and Gauss--Jordan elimination \\
3  & Matrix operations; matrix multiplication as composition \\
4  & Inverse matrices; elementary matrices; LU factorisation \\
5  & Cost of elimination; pivoting and numerical stability \\
6  & Vector spaces and subspaces; fields beyond $\mathbb{R}$ (e.g.\ $\mathbb{F}_2$) \\
7  & Span, linear independence, basis, dimension \\
8  & The four fundamental subspaces; rank; rank--nullity theorem \\
9  & Linear maps: kernel, image, injectivity and surjectivity \\
10 & Matrix of a linear map; change of basis; similarity \\
11 & Determinants: definition, properties, cofactor expansion \\
12 & Determinants: volume, invertibility, Cramer's rule \\
13 & Inner product spaces; norms; Cauchy--Schwarz inequality \\
14 & Orthogonality; orthogonal complements; orthogonal projections \\
15 & Gram--Schmidt process; QR factorisation \\
16 & Least squares and linear regression \\
17 & Eigenvalues and eigenvectors; characteristic polynomial \\
18 & Diagonalisation; powers of matrices; linear recurrences (Fibonacci) \\
19 & Markov chains and PageRank \\
20 & Symmetric matrices; the spectral theorem \\
21 & Quadratic forms; positive definite matrices \\
22 & Singular value decomposition \\
23 & Low-rank approximation; principal component analysis \\
24 & Iterative methods: power iteration; introduction to Jordan form \\
25 & Applications: computer graphics transforms, linear codes over $\mathbb{F}_2$ \\
\bottomrule
\end{tabularx}

\section{Teaching and Learning Methods}
The module is self-paced. All materials are provided on Moodle: lecture notes and
recorded lectures for each unit, tutorial problem sheets with worked solutions,
short quizzes, and the briefs for the two projects. Problem sheets are for practice
and are not marked. Questions are discussed in the module channel on Discord and in
office hours. Independent study includes reading, problem sheets and project work.

\begin{tabularx}{0.6\textwidth}{X r}
\toprule
\textbf{Activity} & \textbf{Hours} \\
\midrule
Lectures            & 50 \\
Labs / Tutorials    & 40 \\
Independent study   & 120 \\
\midrule
\textbf{Total}      & \textbf{210} \\
\bottomrule
\end{tabularx}

\section{Assessment}
Knowledge is assessed with weekly unit quizzes, two programming projects
and a final exam. Labs are used for the programming work.
\begin{itemize}
  \item \textbf{Project 1 -- Solving linear systems:} implement Gaussian elimination
        with partial pivoting and LU factorisation, measure cost and numerical error,
        and fit a model to data by least squares.
  \item \textbf{Project 2 -- Eigenvectors and the SVD:} rank web pages with
        PageRank by power iteration, and compress images with a truncated SVD,
        explaining the trade-off between rank and error.
\end{itemize}
The final exam tests the standard computations, the choice of method for a problem
and short proofs about vector spaces and linear maps.

\begin{tabularx}{\textwidth}{X c c c}
\toprule
\textbf{Component} & \textbf{Weight} & \textbf{Due} & \textbf{LOs} \\
\midrule
Unit quizzes (best 10 of 12) & 20\% & Weekly & LO1--LO6 \\
Project 1: Solving linear systems & 20\% & Week 6 & LO1, LO4, LO7 \\
Project 2: Eigenvectors and the SVD & 20\% & Week 12 & LO5--LO7 \\
Final exam (3 hours, written) & 40\% & End of term & LO1--LO6 \\
\midrule
\textbf{Total}      & \textbf{100\%} & & \\
\bottomrule
\end{tabularx}

\section{Reading List}
\textbf{Essential}
\begin{itemize}
  \item G.\ Strang: \textit{Introduction to Linear Algebra}, 6th ed., Wellesley-Cambridge Press, 2023.
  \item J.\ Hefferon: \textit{Linear Algebra}, 4th ed., 2020. Free online at \url{https://hefferon.net/linearalgebra/}.
\end{itemize}
\textbf{Recommended}
\begin{itemize}
  \item S.\ Axler: \textit{Linear Algebra Done Right}, 4th ed., Springer, 2024 (open access).
  \item S.\ Boyd, L.\ Vandenberghe: \textit{Introduction to Applied Linear Algebra: Vectors, Matrices, and Least Squares}, Cambridge University Press, 2018.
  \item L.\ N.\ Trefethen, D.\ Bau: \textit{Numerical Linear Algebra}, SIAM, 1997.
\end{itemize}

\section{Materials and Licence}
Course materials are published under the
\href{https://creativecommons.org/licenses/by-sa/4.0/}{Creative Commons Attribution-ShareAlike 4.0}
licence unless stated otherwise. See \href{https://wosm.academy/licence/}{wosm.academy/licence} for details.

\section{Policies}
\begin{itemize}
  \item \textbf{Attendance:} Not mandatory. The module is self-paced; work through the units at your own pace within the term.
  \item \textbf{Late submissions:} Accepted, but penalised.
  \item \textbf{Academic integrity:} All submitted work must be your own.
        See the \href{https://wosm.academy/code-of-conduct/}{Code of Conduct}.
\end{itemize}

\end{document}
